% ============================================================
%  USPSC-Cap4-Avaliacao_Experimental_4.3-Analise_Comparativa.tex
%  Seção 4.3 — Análise Comparativa
% ============================================================

\section{Análise Comparativa}
\label{sec:analise_comparativa}

A \autoref{tab:comparativa_geral} consolida os resultados das quatro abordagens avaliadas sobre o mesmo conjunto de teste ($n = 373$), permitindo a comparação direta de acurácia, F1-score por classe e F1-score macro.

\begin{table}[!htbp]
	\caption[Comparativo de desempenho entre as abordagens]{Comparativo de desempenho entre as quatro abordagens ($n = 373$). F1-scores expressos por classe (Pos, Neu, Neg) e como média macro (Macro). Melhor resultado por coluna em negrito.}
	\label{tab:comparativa_geral}
	\centering
	\begin{tabular}{lccccc}
		\hline
		\textbf{Abordagem}  & \textbf{Acurácia} & \textbf{F1 Pos} & \textbf{F1 Neu} & \textbf{F1 Neg} & \textbf{F1 Macro} \\
		\hline
		OpLexicon v3.0      & 0,26              & 0,29            & 0,22            & 0,28            & 0,26              \\
		\hline
		SentiLex-PT         & 0,29              & 0,29            & 0,31            & 0,27            & 0,29              \\
		\hline
		FinBERT-PT-BR       & 0,40              & 0,44            & 0,45            & 0,32            & 0,40              \\
		\hline
		BERTimbau (FT)      & \textbf{0,90}     & \textbf{0,82}   & \textbf{0,94}   & \textbf{0,73}   & \textbf{0,83}     \\
		\hline
	\end{tabular}
	\legend{Fonte: Elaborada pelo autor. (FT) = \textit{fine-tuned} sobre o corpus local.}
\end{table}


\subsection{Hierarquia de Desempenho}
\label{subsec:hierarquia}

Os resultados estabelecem uma hierarquia clara entre as abordagens, ordenada por complexidade metodológica crescente: OpLexicon < SentiLex-PT < FinBERT-PT-BR $\ll$ BERTimbau ajustado. O salto de desempenho mais expressivo ocorre entre o FinBERT-PT-BR e o BERTimbau ajustado: a adaptação ao gênero e ao subdomínio via \textit{fine-tuning} local produziu ganhos de 49,1 pontos percentuais em acurácia e 42,4 pontos em F1-score macro, valores que colocam o BERTimbau ajustado como a única abordagem com desempenho praticamente utilizável para classificação de sentimento neste corpus.

Entre as abordagens léxicas, o SentiLex-PT supera marginalmente o OpLexicon em F1-score macro (29,00\% versus 26,12\%), diferença atribuível, em parte, à maior cobertura morfológica do SentiLex-PT, que inclui entradas anotadas com classes gramaticais e permite maior precisão na recuperação de termos flexionados. Nenhuma das duas abordagens, contudo, supera o limiar de desempenho de um classificador por maioria de classe --- que obteria acurácia de 71,3\% ao classificar todas as instâncias como neutro --- evidenciando que a tarefa não pode ser resolvida por correspondência léxica simples neste tipo de corpus.

\subsection{Desempenho nas Classes Minoritárias}
\label{subsec:classes_minoritarias}

O desempenho nas classes minoritárias (positivo e negativo) é o discriminador mais revelador entre as abordagens. A \autoref{tab:comparativa_minoritarias} detalha precisão e revocação das três classes para cada abordagem.

\begin{table}[!htbp]
	\caption[Precisão e revocação por classe]{Precisão (P) e revocação (R) por classe para cada abordagem}
	\label{tab:comparativa_minoritarias}
	\centering
	\begin{tabular}{lcccccc}
		\hline
		& \multicolumn{2}{c}{\textbf{Positivo}} & \multicolumn{2}{c}{\textbf{Neutro}} & \multicolumn{2}{c}{\textbf{Negativo}} \\
		\textbf{Abordagem} & \textbf{P}   & \textbf{R}              & \textbf{P}  & \textbf{R}            & \textbf{P}   & \textbf{R}              \\
		\hline
		OpLexicon v3.0     & 0,18         & 0,68                    & 0,81        & 0,13                  & 0,20         & 0,46                    \\
		\hline
		SentiLex-PT        & 0,19         & 0,59                    & 0,78        & 0,20                  & 0,18         & 0,48                    \\
		\hline
		FinBERT-PT-BR      & 0,41         & 0,47                    & 0,88        & 0,30                  & 0,20         & 0,88                    \\
		\hline
		BERTimbau (FT)     & 0,84         & 0,80                    & 0,91        & 0,97                  & 0,88         & 0,62                    \\
		\hline
	\end{tabular}
	\legend{Fonte: Elaborada pelo autor.}
\end{table}

As abordagens léxicas apresentam revocação aparentemente elevada para a classe positivo (OpLexicon: 68\%; SentiLex-PT: 59\%), mas com precisão inferior a 20\%, o que indica que a maioria das predições positivas é incorreta --- a proporção de verdadeiros positivos no total de predições positivas é baixíssima. Esse padrão reflete a tendência das abordagens léxicas de superclassificar instâncias como positivo quando o vocabulário não apresenta polaridade negativa explícita, mesmo que o conteúdo seja factual e neutro.

O FinBERT-PT-BR apresenta o inverso para a classe negativo: revocação de 88\% com precisão de apenas 20\%, confirmando o viés pronunciado discutido na \autoref{subsec:res_finbert}. O modelo recupera quase todas as instâncias verdadeiramente negativas, mas ao custo de classificar incorretamente grande parte das neutras como negativas.

O BERTimbau ajustado é a única abordagem que simultaneamente mantém precisão e revocação acima de 60\% para todas as três classes, incluindo as minoritárias. Esse resultado indica que o \textit{fine-tuning} local capacitou o modelo a discriminar os padrões linguísticos específicos das três classes no contexto de publicações financeiras do X/Twitter em português, sem privilegiar sistematicamente nenhuma das classes ao custo das demais.


\subsection{Viés de Classe Neutra}
\label{subsec:vies_neutro}

O tratamento da classe neutra é um indicador relevante da qualidade de cada abordagem para a tarefa específica deste trabalho. Publicações jornalísticas financeiras são predominantemente neutras --- 71,3\% do corpus de teste --- e a correta identificação dessa classe é condição para resultados de classificação de sentimento fiéis à distribuição real de polaridade do corpus. A \autoref{tab:comparativa_geral} mostra que somente o BERTimbau ajustado alcança F1-score acima de 0,80 para a classe neutra (F1 = 0,94). As demais abordagens variam entre 0,22 (OpLexicon) e 0,45 (FinBERT-PT-BR) para essa classe, o que implica que essas abordagens subestimariam sistematicamente a proporção de publicações neutras e exagerariam a fração de publicações com polaridade definida.

